\subsection{婚姻内的同意：结婚证不是"永久同意书"}

婚姻关系中的性同意是一个常被回避的话题：不少人默认"领了结婚证，性就是配偶的权利"。本节纠正这一认知——\textbf{同意是针对每一次具体性行为的，婚姻不产生概括性、永久性的性同意}。

\vspace{0.5em}
\noindent\textbf{核心认知}：

\begin{itemize}
	\item \textbf{"夫妻义务"不是法律义务}：我国法律并未规定配偶必须配合性行为；拒绝性不构成对婚姻的"违约"，更不能用"冷暴力""不尽义务"来指责不同意的伴侣
	\item \textbf{婚姻非正常存续期间的强迫性行为可构成犯罪}：司法实践中，在离婚诉讼期间、分居等婚姻关系已实际破裂的情形下，丈夫强行与妻子发生性行为，已有按强奸罪定罪的判例；此外，婚内强迫性行为还可能构成虐待罪、故意伤害罪
	\item \textbf{《反家庭暴力法》将性暴力纳入家暴范畴}：配偶间的强迫性行为属于家庭暴力，受害人可依法申请告诫书、人身安全保护令（参见"性暴力、家庭暴力"相关章节）
	\item \textbf{同意可以撤回}：即使婚姻存续、即使曾经同意，任何一方都可以在性行为进行中改变主意——此时继续即是强迫
\end{itemize}

\vspace{0.5em}
\noindent\textbf{为什么受害者更难开口}：婚内性强迫的特殊困境在于——受害者常被"夫妻哪有强奸"的观念困住，担心无人相信；经济依附、孩子抚养、家庭名誉等现实牵绊也让人选择沉默；长期处于"不得不从"的状态会造成持续的性厌恶、抑郁与躯体化症状。如果你正处于这样的关系中：这不是你的错，也不是"婚姻该有的样子"，可以拨打 12356 心理援助热线、向妇联求助，或依据《反家庭暴力法》报警（求助资源详见本章"哀伤、丧亲与心理支持资源"及"性暴力"相关章节）。

\vspace{0.5em}
\noindent\textbf{健康婚姻里的性协商}：长期的婚姻里，双方性欲自然会有落差与低谷（参见"性爱意愿不一致与无性婚姻"一节）——健康的路径是沟通与协商，而不是默许或强迫。把"我最近不太想"说出口，比勉强配合一次更保护这段关系。

\begin{tcolorbox}[colback=pink!5!white,colframe=red!50!black,title=贴心小叮咛]
婚姻的意义是两个人自愿的亲密，而不是一方对另一方身体的持有权。\textbf{缔结了婚姻，不等于缔结了对身体的无期限授权}。
\end{tcolorbox}
